\section{International taxation: majority support for taxing the rich to finance LICs}

%The previous section showed that foreign aid can attract majority support when citizens perceive it as effective. %, well targeted, and supported by equitable burden-sharing. AF: deletion
The analysis therefore turns to whether citizens in high-income countries accept %more demanding forms of global redistribution, particularly 
international taxes whose burden would fall predominantly on the richest individuals. % with high income or substantial wealth. AF: changed
The survey evidence presented by \citet{fabre_public_2026} provides a direct test of public acceptance of these instruments. The study combines nationally representative samples of more than 12,000 adults across eleven high-income countries -- the United States, Japan, Russia, Saudi Arabia, France, Germany, Italy, Spain, the United Kingdom, Poland, and Switzerland -- with several survey experiments. % It examines the level and robustness of support across policy designs and country-coverage scenarios, while investigating why substantial stated support has not translated into greater political prominence. % AF: deletion

% AF: ce paragraphe n'a pas sa place ici, c'est pour la dernière section
The gap between stated support and political prominence connects the evidence on global taxation to the political-science literature on issue salience. Voters assign unequal attention and electoral weight to different policy issues. Electoral choices are instead shaped by the issues that become salient and by the positions candidates adopt on those issues \citep{repass_issue_1971,krosnick_role_1988,edwards_explaining_1995,dennison_review_2019}. Consequently, stated support for global redistribution may generate little electoral pressure unless political competition makes the issue sufficiently salient.

\paragraph{International taxation of top incomes}
\citet{fabre_public_2026} first examines support for taxing top incomes globally to finance poverty reduction in low-income countries. Respondents evaluate % AF: ils évaluent une des eux choisie aléatoirement
two variants that differ in the size of the taxable group and the progressivity of the rate schedule. The first imposes an additional 15\% tax on annual after-tax income above \$120,000 at purchasing power parity, thereby targeting the global top 1\%. The second extends taxation to the global top 3\% and applies marginal rates of 15\%, 30\%, and 45\% above the respective thresholds of \$80,000, \$120,000, and \$1 million. The variants are calibrated to close poverty gaps relative to monthly income floors of \$250 and \$400, respectively, and would generate international transfers equivalent to approximately 2\% and 5\% of nominal world income.

%Even these relatively ambitious proposals attract either absolute-majority support or support that substantially exceeds opposition. % AF: deletion
Support reaches 56\% for the top-1\% tax and 50\% for the top-3\% tax, compared with opposition rates of 25\% and 28\%, respectively. Although support falls as the tax base expands, both variants attract substantially more support than opposition, indicating that respondents do not generally reject large-scale redistribution toward low-income countries \citep{fabre_public_2026}. % Say there is majority support among affected respondents overall, though not in all countries for the top 3% variant (Figure S41)

% AF: maybe skip the Zucman tax, this it's not internationally redistributive (talk about it in the paragraph about plausible policies)
The study also examines several internationally redistributive wealth taxes, including a 2\% tax on wealth above \$1 million and a distinct billionaire wealth tax corresponding %more closely 
to the proposal of %associated with  % AF: deletion
Gabriel Zucman. % cite his G20 report
Under the high-income-country coverage scenario, % AF: this coverage is not defined, rephrase
70\% accept a 2\% tax on wealth above \$1 million that allocates 30\% of its revenue to public services in low-income countries. Relative to this scenario, global coverage raises acceptance by 5 % 4.8  AF: changed
percentage points, while narrow %er 
international coverage of just a few progressive countries 
lowers it by a statistically insignificant 1 %1.4 p
ercentage point. These comparisons indicate that acceptance is only slightly affected by %not highly sensitive to  AF: change
incomplete country participation. %Acceptance is highest for the 2\% billionaire wealth tax, at 81\%, compared with 69\% and 64\% for taxes on the global top 1\% and top 3\%, respectively \citep{fabre_public_2026}. % AF: deletion

% AF: ça doit aller en section 8, et il faut aussi mentionner les conjoint experiments du papier NHB
The study further examines whether global redistribution affects choices between hypothetical political programs. In a conjoint experiment, an international millionaire tax allocating 30\% of its revenue to healthcare and education in low-income countries increases the likelihood that a program is preferred %a program's probability of selection AF: changed
by 5 percentage points. By contrast, proposing cuts to development aid reduces it %a program's probability of selection 
by 3 percentage points. Global redistribution thus affects stated program choice, despite its limited prominence in actual electoral competition \citep{fabre_public_2026}. 

% AF: ça doit aller dans section 7
The paper further examines whether respondents are willing to allocate their own tax payments to global objectives. Overall, 41\% agree that their taxes should contribute to solving global problems, compared with 28\% who disagree. %; agreement reaches 60\% among respondents expressing a view. AF: deletion
If the party to which they feel closest joined a global movement combining climate action, millionaire taxation, and poverty reduction in low-income countries, 36\% report that they would become more likely to vote for it, compared with 17\% who would become less likely. Among non-voters who identify with a left-wing party, 46\% report that the movement would make them more likely to vote, whereas 10\% report the opposite. Although hypothetical, these responses indicate that global redistributive taxation may mobilize citizens already ideologically aligned with redistributive and climate-oriented platforms \citep{fabre_public_2026}.

% AF: il s'agit de "Plausible policies" et ça doit venir en début de section
\paragraph{Foreign aid, tax cooperation, and climate-related packages}
Substantial acceptance is not confined to taxes on personal income and wealth. Acceptance reaches 70\% for both debt relief for vulnerable countries and the % a foreign-aid 
commitment 
to allocate % equal to 
0.7\% of GDP % AF: changed
to foreign aid. Acceptance rises to 79\% for the Bridgetown Initiative, %presented as; cite it
a structural reform of international development and climate 
finance. A proposal to raise the global minimum corporate-tax rate from 15\% to 35\% receives 68\% acceptance. Public acceptance therefore encompasses not only explicit redistributive instruments but also institutional reforms of development finance and international tax cooperation \citep{fabre_public_2026}.

% ça doit aller dans section 7
Support also remains substantial when international redistribution forms part of a broader climate-policy package: 68\% prefer a cooperative scenario combining millionaire taxation, vehicle electrification, doubling prices of heating gas, beef and flights, %higher prices for carbon-intensive consumption,  AF: changed
and warming limited to +2\textdegree{}C over a status quo leading to +3\textdegree{}C warming by 2100 \citep{fabre_public_2026}. The result demonstrates substantial acceptance of a costly package jointly addressing climate mitigation, international cooperation, and global redistribution, although it does not identify which component drives support.

% Overall, \citet{fabre_public_2026} finds substantial support for international taxes on top incomes and wealth when revenues are allocated to poverty reduction, public services, or climate action in low-income countries. AF: changed
Overall, \citet{fabre_public_2026} finds substantial support for international transfers when financed by  taxes on top incomes and wealth. 
% AF: pas d'accord avec cette conclusion: taxer les riches pour réduire la pauvreté à l'échelle mondiale, c'est précisément ça le redistribution mondiale. C'est les résultats sur foreign aid qui montre un soutien conditionnel (à l'absence de corruption)
% The findings nevertheless fall short of demonstrating broad and unconditional enthusiasm for global redistribution. The evidence instead indicates conditional acceptance, with support concentrated on taxes targeting very wealthy households and linked to international cooperation, poverty reduction, and climate action. 
% The central puzzle is thus the persistence of political under-provision and low issue salience despite majority or plurality support for several international redistributive policies.

